%------------------------------------------------------------------------------%

\usepackage{calculo_varias_variaveis-1}

\title{Propriedades do Vetor Gradiente}

%------------------------------------------------------------------------------%
\begin{document}

\addtitle

%------------------------------------------------------------------------------%
\section{Vetor Gradiente}

%------------------------------------------------------------------------------%
\begin{frame}[t]{Vetor Gradiente}
  Se uma função \(f(x, y)\) possui derivadas parciais,
  seu \emph{vetor gradiente} é

  \vspace{0.99\baselineskip}
  \[
    \nabla f
    \;=\; \D{f}{x}\veci + \D{f}{y}\vecj
    \;=\; \left(\begin{array}{c}
            \displaystyle\D{f}{x} \\[5mm]
            \displaystyle\D{f}{y}
          \end{array}\right)
  \]
\end{frame}

%------------------------------------------------------------------------------%
\begin{frame}[t]{Vetor Gradiente 3D}
  Se uma função \(f(x, y, z)\) possui derivadas parciais,
  seu \emph{vetor gradiente} é

  \[
    \nabla f
    \;=\; \D{f}{x}\veci + \D{f}{y}\vecj + \D{f}{z}\veck
    \;=\; \left(\begin{array}{c}
            \displaystyle\D{f}{x} \\[5mm]
            \displaystyle\D{f}{y} \\[5mm]
            \displaystyle\D{f}{z}
          \end{array}\right)
  \]
\end{frame}

%------------------------------------------------------------------------------%
\section{Gradientes e Crescimento da Função}

%------------------------------------------------------------------------------%
\begin{frame}{Derivada Direcional}
  \begin{align*}
    \onslide<+->{D_u f}
    \onslide<+->{& = \nabla f \cdot u                        \\[5mm]}
    \onslide<+->{& = \norm{\nabla f}\; \norm{u} \cos(\theta)     \\[5mm]}
    \onslide<+->{& = \norm{\nabla f} \cos(\theta)}
  \end{align*}
\end{frame}

%------------------------------------------------------------------------------%
\begin{frame}{Derivada Direcional}
\begin{enumerate}
  \setlength\itemsep{1.5em}
  \item A função aumenta mais rapidamente quando
        \pause
        \tabto{82mm}\emph{$\cos(\theta)=1$},
        \pause
        \emph{$\hphantom{-}\quad\theta=0$} \\[2mm]
        \pause
        Direção do gradiente
        \pause
  \item A função descresse mais rapidamente quando
        \pause
        \tabto{82mm}\emph{$\cos(\theta)=-1$},
        \pause
        \emph{$\quad\theta=\pi$} \\[2mm]
        \pause
        Direção oposta ao gradiente
        \pause
  \item Se $u$ for \emph{ortogonal} a \(\nabla f\neq 0\)
        \pause
        ele aponta para uma direção de \emph{variação zero}
\end{enumerate}
\end{frame}

%------------------------------------------------------------------------------%
\begin{frame}{Gradientes e Crescimento da Função}
  \begin{enumerate}[<+->]
    \setlength\itemsep{1.5em}
    \item Direção de maior crescimento da função $f$ no ponto \({x, y}\)
          \[
            v_M = \nabla f(x, y)
          \]
    \item Direção de maior decrescimento da função $f$ no ponto \({x, y}\)
          \[
            v_m = -\nabla f(x, y)
          \]
    \item Direção de variação nula da função $f$ no ponto \({x, y}\)
          \[
            \nabla f(x, y) \cdot v_o  = 0
          \]
  \end{enumerate}
\end{frame}

%------------------------------------------------------------------------------%
\addtocounter{exemplo}{1}
\begin{frame}{Exemplo \theexemplo}
  Analisando a função \(f(x, y) = x^2 + \frac{y^2}{2}\)
  no ponto \((1, 1)\)

  Encontre as direções nas quais:

  a)\; $f$ cresce mais rapidamente

  b)\; $f$ decresce mais rapidamente

  c)\; $f$ não varia
\end{frame}

\begin{frame}{Exemplo \theexemplo}
  Calculando as derivadas parciais

  \[
    \pause
    \D{f}{x}
    \pause
    \:=\; \D{}{x}\left(x^2 + \frac{y^2}{2}\right)
    \pause
    \:=\; 2x
  \]

  \[
    \pause
    \D{f}{y}
    \pause
    \:=\; \D{}{y}\left(x^2 + \frac{y^2}{2}\right)
    \pause
    \:=\; y
  \]
\end{frame}

\begin{frame}{Exemplo \theexemplo}
  Montando o gradiente
  \pause
  \[
    \nabla f
    \pause
    = \Vetor{\D{f}{x}}{\D{f}{y}}
    \pause
    = \Vetor{\D{}{x}\left(x^2 + \frac{y^2}{2}\right)}
            {\D{}{y}\left(x^2 + \frac{y^2}{2}\right)}
    \pause
    = \vetor{2x}{y}
    \]
    \pause
    Avaliando o gradiente no ponto solicitado
    \[
    \pause
    \nabla f(1, 1)
    \pause
    \;=\; \vetor{2x}{y}\evalat{(1, 1)}{}
    \pause
    \;=\; \vetor{2}{1}
  \]
\end{frame}

\begin{frame}{Exemplo \theexemplo}
  a)\; $f$ cresce mais rapidamente
  \pause
  na direção do gradiente
  \[
    \pause
    v_M
    \pause
    = \nabla f(1, 1)
    \pause
    = \vetor{2}{1}
  \]

  \pause
  b)\; $f$ decresce mais rapidamente
  \pause
  na direção oposta ao gradiente
  \[
    \pause
    v_m
    \pause
    = -\nabla f (1, 1)
    \pause
    = \vetor{-2}{-1}
  \]
\end{frame}

\begin{frame}{Exemplo \theexemplo}
  \onslide<+->{c)\; $f$ não varia}
  \onslide<+->{na direção, $v_o$, ortogonal ao gradiente}
  \begin{align*}
    \onslide<+->{\nabla f(1, 1) \cdot v_o              & = 0 \\[2mm]}
    \onslide<+->{\vetor{2}{1} \cdot \vetor{v_1}{v_2}   & = 0 \\[2mm]}
    \onslide<+->{2v_1 + v_2                            & = 0 \\[2mm]}
    \onslide<+->{v_2                                   & = -2v_1}
  \end{align*}

  \onslide<+->{As direções são
  \(
    v_o = \vetor{1}{-2}
  \)}
  \onslide<+->{e
  \(
    v_o = \vetor{-1}{2}
  \)}
\end{frame}

%------------------------------------------------------------------------------%
%\addtocounter{exemplo}{1}
%\begin{frame}{Exemplo \theexemplo}
%\end{frame}
%
%\begin{frame}{Exemplo \theexemplo}
%\end{frame}

%------------------------------------------------------------------------------%
\section{Gradientes e Curvas de Nível}

%------------------------------------------------------------------------------%
\begin{frame}{Curvas de Nível}
  Curva de Nível é o conjunto de pontos \((x, y)\) onde
  \[
    f(x, y) = c
  \]
  para um valor constante $c$

  \pause
  \vspace{0.5\baselineskip}
  Em alguns casos, é possível escrever essa região como uma curva paramétrica
  \[
    x = g(t)
    \qquad
    y = h(t)
  \]
  \pause
  Vetorialmente temos
  \[
    r(t)
    = g(t)\veci + h(t)\vecj
    = \vetor{g(t)}{h(t)}
  \]
\end{frame}

%------------------------------------------------------------------------------%
\begin{frame}{Tangente de uma Curva de Nível}
\begin{columns}
\column{0.5\linewidth}
  \begin{align*}
    \onslide<+->{f\big(g(t), h(t)\big) & = c \\[2mm]}
    \onslide<+->{\frac{d}{dt} f\big(g(t), h(t)\big) & = 0 \\[2mm]}
    \onslide<+->{\D{f}{x}\frac{dg}{dt} + \D{f}{y}\frac{dh}{dt} & = 0  \\[2mm]}
    \onslide<+->{\Vetor{\D{f}{x}}{\D{f}{y}} \cdot
                 \Vetor{\frac{dg}{dt}}{\frac{dh}{dt}} & = 0}
  \end{align*}
\column{0.5\linewidth}
  \onslide<+->{\[\nabla f \cdot \frac{dr}{dt} = 0\]}

  \onslide<+->{O gradiente é ortogonal ao vetor tangente à curva de nível}
\end{columns}
\end{frame}

%------------------------------------------------------------------------------%
\begin{frame}{Funções Diferenciáveis}
  Em todo ponto \((x_0, y_0)\) no domínio de uma \emph{função diferenciável}
  \(f(x, y)\),

  \pause
  o \emph{gradiente} de $f$
  \pause
  é normal à \emph{curva de nível} que passa por \((x_0, y_0)\)
\end{frame}

\framedgraphic{Gradientes e Curvas de Nível}{gradiente_curva_nivel}{}

%------------------------------------------------------------------------------%
\addtocounter{exemplo}{1}
\begin{frame}{Exemplo \theexemplo}
  Dada a função
  \(
    f(x, y) = (x-2)^2 + (y-1)^2
  \)

  a) Determine e esboce a curva de nível \(f(x, y) = 4\)

  b) Calcule e esboce o gradiente de $f$ nos pontos
  \((0,  1)\),
  \((2, -1)\),
  \((2,  3)\) e
  \((4,  1)\)
\end{frame}

%------------------------------------------------------------------------------%
\begin{frame}{Exemplo \theexemplo}
  \onslide<+->{a) Determinando a curva de nível}
  \begin{align*}
    \onslide<+->{f(x, y)           & = 4}   \\
    \onslide<+->{(x-2)^2 + (y-1)^2 & = 2^2}
  \end{align*}
  \onslide<+->{Equação da circunferência centrada em \((2, 1)\) e raio 2}
\end{frame}

%------------------------------------------------------------------------------%
\begin{frame}{Exemplo \theexemplo}
  b) Calculando as derivadas parciais de $f$

  \pause
  \[
    \D{f}{x}
    \pause
    \;=\; \D{}{x}\left[(x-2)^2 + (y-1)^2\right]
    \pause
    \;=\; 2(x-2)\D{}{x}(x-2)
    \pause
    \;=\; 2(x-2)
  \]

  \pause
  \[
    \D{f}{y}
    \pause
    \;=\; \D{}{y}\left[(x-2)^2 + (y-1)^2\right]
    \pause
    \;=\; 2(y-1)\D{}{y}(y-1)
    \pause
    \;=\; 2(y-1)
  \]

  \pause
  Gradiente de $f$
  \pause
  \[
    \nabla f(x, y)
    \pause
    \;=\; \nabla\left[(x-2)^2 + (y-1)^2\right]
    \pause
    \;=\; \begin{pmatrix}
      2(x-2) \\[2mm] 2(y-1)
    \end{pmatrix}
  \]
\end{frame}

%------------------------------------------------------------------------------%
\begin{frame}{Exemplo \theexemplo}
  Calculando o gradiente nos pontos
  \((0,  1)\) e
  \((2, -1)\)

  \pause
  \[
    \nabla f(0, 1)
    \pause
    \;=\; \begin{pmatrix}
      2(x-2) \\[2mm] 2(y-1)
    \end{pmatrix} \evalat{(0, 1)}{}
    \pause
    \;=\; \begin{pmatrix}
    2(0-2) \\[2mm] 2(1-1)
    \end{pmatrix}
    \pause
    \;=\; \begin{pmatrix}
    -4 \\[2mm] 0
    \end{pmatrix}
  \]

  \pause
  \[
    \nabla f(2, -1)
    \pause
    \;=\; \begin{pmatrix}
      2(x-2) \\[2mm] 2(y-1)
    \end{pmatrix} \evalat{(2, -1)}{}
    \pause
    \;=\; \begin{pmatrix}
      2(2-2) \\[2mm] 2(-1-1)
    \end{pmatrix}
    \pause
    \;=\; \begin{pmatrix}
      0 \\[2mm] -4
    \end{pmatrix}
  \]
\end{frame}

%------------------------------------------------------------------------------%
\begin{frame}{Exemplo \theexemplo}
  Calculando o gradiente nos pontos
  \((2,  3)\) e
  \((4,  1)\)

  \pause
  \[
    \nabla f(2, 3)
    \pause
    \;=\; \begin{pmatrix}
      2(x-2) \\[2mm] 2(y-1)
    \end{pmatrix} \evalat{(2, 3)}{}
    \pause
    \;=\; \begin{pmatrix}
      2(2-2) \\[2mm] 2(2-1)
    \end{pmatrix}
    \pause
    \;=\; \begin{pmatrix}
      0 \\[2mm] 4
    \end{pmatrix}
  \]

  \pause
  \[
    \nabla f(4, 1)
    \pause
    \;=\; \begin{pmatrix}
      2(x-2) \\[2mm] 2(y-1)
    \end{pmatrix} \evalat{(4, 1)}{}
    \pause
    \;=\; \begin{pmatrix}
      2(4-2) \\[2mm] 2(1-1)
    \end{pmatrix}
    \pause
    \;=\; \begin{pmatrix}
      4 \\[2mm] 0
    \end{pmatrix}
  \]
\end{frame}

%------------------------------------------------------------------------------%
\addtocounter{exemplo}{1}
\begin{frame}[t]{Exemplo \theexemplo}
  Encontre uma equação para a tangente à elipse
  \[
    \frac{x^2}{4} + y^2 = 2
  \]
  no ponto \((-2, 1)\)
\end{frame}

\begin{frame}[t]{Exemplo \theexemplo}
  Encontre uma equação para a tangente à elipse
  \[
    \frac{x^2}{4} + y^2 = 2
  \]
  no ponto \((-2, 1)\)

  \addgraphic{9}{6.5,2.5}{tangente_elipse}
\end{frame}

\begin{frame}{Exemplo \theexemplo}
  Precisamos notar que a elipse é uma curva de nível da função
  \[
    f(x, y) = \frac{x^2}{4} + y^2
  \]

  \pause
  Calculando as derivadas parciais de $f$ temos
  \[
    \pause
    \D{f}{x}
    \pause
    = \D{}{x}\left(\frac{x^2}{4} + y^2\right)
    \pause
    = \frac{x}{2}
  \]
  \[
    \pause
    \D{f}{y}
    \pause
    = \D{}{y}\left(\frac{x^2}{4} + y^2\right)
    \pause
    = 2y
  \]
\end{frame}

\begin{frame}{Exemplo \theexemplo}
  Então o gradiente de $f$, no ponto \((-2, 1)\), é
  \pause
  \[
    \nabla f(-2, 1)
    \pause
    \;=\; \Vetor{\frac{x}{2}}{2y}\evalat{(-2, 1)}{}
    \pause
    = \vetor{-1}{2}
  \]
\end{frame}

\begin{frame}{Exemplo \theexemplo}
  \onslide<+->{A reta tangente é ortogonal ao vetor gradiente}
  \begin{columns}
  \column{0.5\linewidth}
    \begin{align*}
      \onslide<+->{\nabla f(x_0, y_0) \cdot \vetor{x-x_0}{y-y_0} & = 0 \\[3mm]}
      \onslide<+->{\nabla f(-2, 1)    \cdot \vetor{x-(-2)}{y-1} & = 0 \\[3mm]}
      \onslide<+->{\vetor{-1}{2} \cdot \vetor{x+2}{y-1} & = 0}
    \end{align*}
  \column{0.5\linewidth}
  \begin{align*}
    \onslide<+->{-(x+2) + 2(y-1) & = 0 \\}
    \onslide<+->{-x +2y -4 & = 0 \\}
    \onslide<+->{-x +2y    & = 4 \\}
    \onslide<+->{ x        & = 2y-4}
  \end{align*}
  \end{columns}
\end{frame}

%------------------------------------------------------------------------------%
\section{Propriedades Algébricas}

%------------------------------------------------------------------------------%
\begin{frame}{Propriedades Algébricas do Gradiente}
  \begin{enumerate}[<+->]
    \setlength\itemsep{1.5em}
    \item Soma                        \tabto{55mm} \(\nabla(f+g) = \nabla f + \nabla g\)
    \item Diferença                   \tabto{55mm} \(\nabla(f-g) = \nabla f - \nabla g\)
    \item Multiplicação por constante \tabto{55mm} \(\nabla(kf) = k\nabla f\)
    \item Produto                     \tabto{55mm} \(\nabla(fg) = f\nabla g + g\nabla f\)
    \item Quociente                   \tabto{55mm}
      \(\nabla\left(\frac{f}{g}\right) = \frac{g\nabla f - f\nabla g}{g^2}\)
  \end{enumerate}
\end{frame}

%------------------------------------------------------------------------------%
\addtocounter{exemplo}{1}
\begin{frame}{Exemplo \theexemplo}
  Dadas \; \(f(x, y)=x^2-2y\) \; e \; \(g(x, y)=xy^2\)

  a)\; calcule os gradientes de $f$ e $g$

  b)\; calcule \(\nabla(f-g)\)

  c)\; calcule \(\nabla(fg)\)
\end{frame}

\begin{frame}{Exemplo \theexemplo}
  \[
    \nabla f
    \pause
    \;=\; \Vetor{\D{f}{x}}{\D{f}{y}}
    \pause
    \;=\; \Vetor{\D{}{x}\left(x^2-2y\right)}{\D{}{y}\left(x^2-2y\right)}
    \pause
    \;=\; \vetor{2x}{-2}
  \]

  \pause
  \[
    \nabla g
    \pause
    \;=\; \Vetor{\D{g}{x}}{\D{g}{y}}
    \pause
    \;=\; \Vetor{\D{}{x}\left(xy^2\right)}{\D{}{y}\left(xy^2\right)}
    \pause
    \;=\; \vetor{y^2}{2xy}
  \]
\end{frame}

\begin{frame}{Exemplo \theexemplo}
  \begin{align*}
    \onslide<+->{\nabla(f-g)}
    \onslide<+->{& = \nabla f - \nabla g \\[2mm]}
    \onslide<+->{& = \vetor{2x}{-2} - \vetor{y^2}{2xy} \\[2mm]}
    \onslide<+->{& = \vetor{2x-y^2}{-2-2xy}}
  \end{align*}
\end{frame}

\begin{frame}{Exemplo \theexemplo}
  \begin{align*}
    \onslide<+->{\nabla(fg)}
    \onslide<+->{& = f\nabla g + g\nabla f                   \\[2mm]}
    \onslide<+->{& = \big(x^2-2y\big)\vetor{y^2}{2xy}
                   + \big(xy^2  \big)\vetor{2x}{2}           \\[2mm]}
    \onslide<+->{& = \vetor{(x^2-2y)y^2}{(x^2-2y)2xy}
                   + \vetor{(xy^2)2x}{(xy^2)2}               \\[2mm]}
    \onslide<+->{& = \vetor{x^2y^2-2y^3}{2x^3y-4xy^2}
                   + \vetor{2x^2y^2}{2xy^2}                  \\[2mm]}
    \onslide<+->{& = \vetor{3x^2y^2 - 2y^3}{2x^3y - 2xy^2}}
  \end{align*}
\end{frame}

%------------------------------------------------------------------------------%
\section{Lista Mínima}

%------------------------------------------------------------------------------%
\begin{frame}{Lista Mínima}
  Cálculo Vol. 2 do Thomas 12\textsuperscript{a} ed. -- Seção 14.5
  \begin{enumerate}
    \item Estudar o texto da seção
    \item Resolver os exercícios: 20, 22, 24, 26, 28, 29
  \end{enumerate}

  \vfill
  Atenção: A prova é baseada no livro, não nas apresentações
\end{frame}

%------------------------------------------------------------------------------%
\end{document}
%------------------------------------------------------------------------------%
